\documentclass[11pt,a4paper]{article}
\input{preambolo}

\title{\bfseries Simulazione quantistica di un dimero di spin:\\
dallo stato fondamentale al sistema aperto}
\author{Samuele Schiavi \\ \small Tesi di laurea triennale in Fisica, Università degli Studi di Parma \\ \small Relatore: Prof.\ Alessandro Chiesa}
\date{}

\begin{document}
\maketitle
\thispagestyle{fancy}

\begin{abstract}
\noindent
Si studia un dimero di spin $1/2$ accoppiati con interazione di Heisenberg
isotropa, campo Zeeman e un termine antisimmetrico di
Dzyaloshinskii--Moriya (DM), come primo banco di prova per la simulazione
quantistica di sistemi di spin molecolari su hardware attuale. Nella prima
parte il sistema è trattato come chiuso: si determina lo spettro esatto, si
costruisce e si confronta un insieme di ansatz variazionali (VQE) per lo
stato fondamentale, si implementa l'evoluzione temporale reale
$e^{-iHt}$ tramite decomposizione di Suzuki--Trotter e si misura, con un
circuito ad ancilla, un correlatore dinamico a due tempi. Nella seconda
parte lo stesso sistema è trattato come aperto: introdotto un modello di
rumore calibrato su hardware IBM reale (canali depolarizzanti a uno e due
qubit più errore di lettura), si ripercorre l'intera catena --- stato
fondamentale, evoluzione temporale, correlatori --- misurando la
degradazione indotta dal rumore e individuando, per la dinamica, un numero
ottimale di passi di Trotter $N^*$ che bilancia l'errore di
discretizzazione con il rumore di gate accumulato. Si esamina infine la
dipendenza di $N^*$ dai parametri del canale di rumore e due estensioni
mirate: l'effetto di un'ottimizzazione variazionale che veda il rumore
durante la ricerca dei parametri, e l'effetto di un errore di lettura
asimmetrico.
\end{abstract}

\tableofcontents

\bigskip

\section{Il sistema fisico}
\label{sec:sistema}

\subsection{Hamiltoniana e corrispondenza spin--qubit}

Il sistema studiato è costituito da due spin $1/2$ accoppiati, descritti
dall'Hamiltoniana
\begin{equation}
H = \underbrace{J\,(X_1X_2+Y_1Y_2+Z_1Z_2)}_{\text{scambio isotropo}}
  + \underbrace{b\,(Z_1+Z_2)}_{\text{campo Zeeman}}
  + \underbrace{D\,(X_1Z_2-Z_1X_2)}_{\text{termine DM}},
\label{eq:H}
\end{equation}
dove $X_i,Y_i,Z_i$ sono le matrici di Pauli sul sito $i=1,2$. I primi due
termini descrivono uno scambio di Heisenberg isotropo e l'accoppiamento
Zeeman con un campo magnetico uniforme lungo $z$; il terzo, antisimmetrico
nello scambio dei siti, è un termine di Dzyaloshinskii--Moriya (DM), che
compare in molecole magnetiche prive di un centro di inversione tra i due
ioni accoppiati. Fissando $J=1$ come unità di energia, il sistema è
descritto da due parametri adimensionali, $b/J$ e $D/J$.

Ogni spin $1/2$ si mappa direttamente su un qubit, senza alcuna
trasformazione fermionica intermedia: è questa corrispondenza diretta,
propria dei sistemi di spin, a rendere le molecole magnetiche un banco di
prova naturale per l'hardware quantistico attuale, a differenza di sistemi
elettronici che richiederebbero un mapping (Jordan--Wigner o simili) e un
numero di qubit ben maggiore del numero di gradi di libertà fisici.
Il modello isotropo con termine DM qui adottato segue l'impostazione di
Crippa \emph{et al.}\ \cite{crippa2021}, dove il caso a due spin è discusso
come primo esempio di simulazione VQE per sistemi di spin molecolari.

\subsection{Spettro esatto e incrocio di livelli}

Con $D=0$, l'Hamiltoniana \eqref{eq:H} conserva la magnetizzazione totale
$M=(Z_1+Z_2)/2$ e commuta con lo spin totale al quadrato $S^2$; il problema
si risolve allora in forma chiusa nella base $\ket{S,M}$, con autovalori
\begin{equation}
E(S,M) = 2J\,S(S{+}1) - 3J + 2bM.
\label{eq:spettro-esatto}
\end{equation}
Lo spazio a quattro dimensioni si decompone in un singoletto ($S=0$,
$M=0$, energia $E=-3J$, indipendente dal campo) e un tripletto ($S=1$), le
cui tre componenti ($M=+1,0,-1$) si separano linearmente in $b$. In termini
degli stati di registro dei due qubit,
\begin{equation}
\ket{1,1}=\ket{00},\qquad
\ket{1,0}=\tfrac{1}{\sqrt2}(\ket{01}+\ket{10}),\qquad
\ket{0,0}=\tfrac{1}{\sqrt2}(\ket{01}-\ket{10}),\qquad
\ket{1,-1}=\ket{11}.
\label{eq:base-registro}
\end{equation}
I due stati con $M=\pm1$ sono separabili (stati di prodotto); i due con
$M=0$ --- il singoletto e la componente centrale del tripletto --- sono
entrambi genuinamente entangled.

\begin{figure}[htbp]
\centering
\includegraphics[width=0.78\textwidth]{figure/fig01_spettro}
\caption{Spettro dei quattro livelli in funzione del campo, a $D=0$. Il
singoletto (energia costante) non risente del campo; le tre componenti del
tripletto si aprono a ventaglio linearmente in $b$. Lo stato fondamentale
(curva più bassa) è il singoletto per $B/J<2$ e la componente $M=-1$ del
tripletto per $B/J>2$; a $B/J=2$ le due energie si incrociano esattamente,
senza deviare.}
\label{fig:spettro}
\end{figure}
\FloatBarrier

A $B/J=2$ le energie del singoletto e della componente $M=-1$ del
tripletto si incrociano esattamente (Fig.~\ref{fig:spettro}): un incrocio
vero, non un avvicinamento, perché i due stati hanno magnetizzazione
diversa e, finché $H$ conserva $M$, l'elemento di matrice che li
collegherebbe è nullo per ragioni di simmetria,
\begin{equation}
\bra{S{=}0,M{=}0}\,H_{J}+H_{b}\,\ket{S{=}1,M{=}{-}1} = 0.
\end{equation}
In prossimità dell'incrocio lo stato fondamentale cambia identità di colpo:
da uno stato completamente entangled (il singoletto) a uno stato di
prodotto ($\ket{\downarrow\downarrow}$). Verificato numericamente per
diagonalizzazione diretta della matrice $4\times4$: lo spettro coincide con
la formula analitica \eqref{eq:spettro-esatto} su tutta la griglia in $b$,
con scarto massimo $8.9\times10^{-16}$ (precisione di macchina).

\subsection{Il ruolo del termine di Dzyaloshinskii--Moriya}

Il termine DM non conserva la magnetizzazione totale: la sua non
commutatività con $S_z$ e $S^2$ è quantificabile tramite la norma di
Frobenius del commutatore,
\begin{equation}
\lVert [H_D, S_z] \rVert = 0.566, \qquad
\lVert [H_D, S^2] \rVert = 1.131 \qquad (D/J=0.2),
\end{equation}
entrambe non nulle: il termine può mescolare stati con $M$ diverso e quindi
aprire l'incrocio di Fig.~\ref{fig:spettro}. L'effetto è visibile su due
osservabili (Fig.~\ref{fig:gap}): il gap tra stato fondamentale e primo
eccitato, che con $D=0$ si annulla esattamente a $B/J=2$ ma con $D/J=0.2$
resta finito (minimo $0.565$, spostato a $B/J\approx2.007$); e la
magnetizzazione dello stato fondamentale, che passa da un salto netto a un
crossover continuo.

\begin{figure}[htbp]
\centering
\includegraphics[width=\textwidth]{figure/fig02_gap_magnetizzazione}
\caption{\textbf{Sinistra:} gap tra stato fondamentale e primo eccitato in
funzione del campo. Con $D=0$ (nero) si annulla a $B/J=2$; con $D/J=0.2$
(arancio) resta finito, minimo $0.565$. \textbf{Destra:} magnetizzazione
dello stato fondamentale; il salto discontinuo diventa un crossover
continuo.}
\label{fig:gap}
\end{figure}
\FloatBarrier

Questo singolo fatto --- la rottura esplicita della conservazione di $M$ ---
è la causa comune di tre conseguenze sviluppate nel seguito:

\begin{enumerate}
\item rende ben posto il problema variazionale su tutto il range di campo,
perché lo stato fondamentale è unico anche dove, senza DM, sarebbe
degenere;
\item rompe la simmetria su cui un ansatz variazionale efficiente sarebbe
naturalmente costruito (Sez.~\ref{sec:vqe});
\item impedisce la fattorizzazione esatta dell'evoluzione temporale
$e^{-iHt}$, rendendo necessaria una decomposizione di Suzuki--Trotter
(Sez.~\ref{sec:trotter}).
\end{enumerate}

Si tratta di una rottura \emph{esplicita} della simmetria --- un termine
aggiunto all'Hamiltoniana --- non di una rottura spontanea: l'Hamiltoniana
stessa, con $D\neq0$, non possiede più quella simmetria.

\section{Stato fondamentale variazionale (VQE)}
\label{sec:vqe}

\subsection{Il metodo}

Il Variational Quantum Eigensolver (VQE) determina un'approssimazione allo
stato fondamentale minimizzando l'energia media su un circuito
parametrico. Si prepara $\ket{\psi(\boldsymbol\theta)}$ con un circuito che
dipende da parametri $\boldsymbol\theta$, si misura
$E(\boldsymbol\theta)=\bra{\psi(\boldsymbol\theta)}H\ket{\psi(\boldsymbol\theta)}$,
e un ottimizzatore classico aggiorna $\boldsymbol\theta$ per abbassare
$E$. Il principio variazionale garantisce $E(\boldsymbol\theta)\ge E_0$
per ogni $\boldsymbol\theta$. Poiché sul dimero l'energia esatta è nota
(Sez.~\ref{sec:sistema}), il VQE qui non serve a scoprirla, ma a misurare
quanto bene un dato circuito riesca a rappresentare lo stato fondamentale
vero: la grandezza di riferimento è la fidelity
$\mathcal F=\lvert\braket{\psi_{\rm VQE}}{\psi_0}\rvert$, o, sui
sottospazi degeneri, la somma delle proiezioni su tutti gli autovettori
degeneri con l'energia più bassa.

\subsection{Due filosofie di ansatz}

Si confrontano due strategie di costruzione del circuito.

\paragraph{Ansatz generico (hardware-efficient, HA).} Rotazioni $R_y$ a un
qubit alternate a porte $CZ$, in uno schema regolare che non incorpora
alcuna informazione sulla fisica del problema: $6$ parametri
(Fig.~\ref{fig:circ-ha}).

\begin{figure}[htbp]
\centering
\includegraphics[width=0.5\textwidth]{figure/circ_ha}
\caption{Ansatz generico hardware-efficient: rotazioni $R_y$ alternate a
porte $CZ$, $6$ parametri.}
\label{fig:circ-ha}
\end{figure}
\FloatBarrier

\paragraph{Ansatz fisicamente motivato (PMA), versione base.} Costruito a
partire da ciò che si sa dello stato fondamentale a $D=0$: sotto l'incrocio
si prepara direttamente il singoletto, sopra si prepara
$\ket{\downarrow\downarrow}$, e si aggiunge un blocco a un solo parametro
--- due CNOT attorno a una rotazione $R_y(\theta)$ (Fig.~\ref{fig:cnotry})
--- che ruota \emph{dentro} il sottospazio a due dimensioni di partenza,
senza mai uscirne: $\{\ket{01},\ket{10}\}$ sotto l'incrocio,
$\{\ket{00},\ket{11}\}$ sopra.

\begin{figure}[htbp]
\centering
\includegraphics[width=0.3\textwidth]{figure/circ_cnotry}
\caption{Blocco a un parametro dell'ansatz PMA base: due CNOT attorno a
un'unica $R_y(\theta)$.}
\label{fig:cnotry}
\end{figure}
\FloatBarrier

Il circuito ha due rami distinti, scelti dalla stessa condizione $b/J\lessgtr2$
che separa i due tratti dello spettro (Fig.~\ref{fig:pma-base}): sotto
l'incrocio prepara il singoletto con le porte $X,H,\mathrm{CNOT},X$, sopra
prepara $\ket{11}$ con due porte $X$; in entrambi i casi segue lo stesso
blocco parametrico di Fig.~\ref{fig:cnotry}. In totale, un solo parametro
libero contro i sei dell'ansatz generico.

\begin{figure}[htbp]
\centering
\begin{minipage}{0.48\textwidth}
\centering
\includegraphics[width=0.95\textwidth]{figure/circ_pma_base}
\vspace{2pt}
\textbf{Ramo $b/J<2$}
\end{minipage}
\hfill
\begin{minipage}{0.48\textwidth}
\centering
\includegraphics[width=0.95\textwidth]{figure/circ_pma_base_ramo2}
\vspace{2pt}
\textbf{Ramo $b/J\geq2$}
\end{minipage}
\caption{Ansatz PMA base nei suoi due rami: preparazione del singoletto
(sinistra) o di $\ket{11}$ (destra), seguita dallo stesso blocco
parametrico di Fig.~\ref{fig:cnotry}.}
\label{fig:pma-base}
\end{figure}
\FloatBarrier

\subsection{Confronto: successo senza DM, fallimento strutturale con DM}

Senza termine DM, i due ansatz sono equivalenti su tutto il range di campo:
fidelity $=1$ per entrambi (Fig.~\ref{fig:confronto}, pannello sinistro).
Imporre la simmetria giusta (conservazione di $M$) riduce il costo di un
fattore sei senza perdita di accuratezza --- il risultato metodologico
centrale di Crippa \emph{et al.}~\cite{crippa2021}.

Con il termine DM acceso, l'ansatz generico resta perfetto su tutto il
range; l'ansatz PMA base crolla a $\mathcal F\approx0.70$ nell'intorno
dell'incrocio, con un minimo di $0.491$ esattamente a $B/J=2$
($D/J=0.2$; Fig.~\ref{fig:confronto}, pannello destro).

\begin{figure}[htbp]
\centering
\includegraphics[width=\textwidth]{figure/fig03_vqe_ha_vs_pma}
\caption{Fidelity rispetto allo stato fondamentale esatto, ansatz generico
(HA) contro ansatz fisicamente motivato base (PMA base). \textbf{Sinistra:}
$D=0$, entrambi perfetti. \textbf{Destra:} $D/J=0.2$, l'ansatz generico
resta perfetto, il PMA base crolla nell'intorno dell'incrocio.}
\label{fig:confronto}
\end{figure}
\FloatBarrier

La causa del fallimento non è il numero di parametri. Vicino all'incrocio,
con il DM acceso, lo stato fondamentale è una sovrapposizione genuina fra
settori di magnetizzazione diversa,
\begin{equation}
\ket{\psi_0} = \alpha\ket{\text{singoletto}} + \beta\ket{\downarrow\downarrow},
\end{equation}
mentre il blocco CNOT--$R_y$--CNOT, per costruzione, non lascia mai il
sottospazio di partenza. Ripetendo il blocco $K=1,2,3,4$ volte, ciascuna
copia con angolo indipendente, la fidelity a $B/J=2$ resta identica a ogni
cifra mostrata ($0.491178$): il circuito è già al meglio possibile dentro
il settore in cui è confinato, e aggiungere parametri dello stesso tipo
non aggiunge direzioni raggiungibili. Il limite è strutturale (la
simmetria imposta), non di capacità espressiva.

\subsection{Riparazione: blocco RBS e ansatz esteso}

Il blocco $M$-conservante usato per la riparazione non è
CNOT--$R_y$--CNOT ma una rotazione di Givens,
\begin{equation}
\mathrm{RBS}(\varphi) =
\begin{pmatrix}
1 & 0 & 0 & 0\\
0 & \cos\varphi & -\sin\varphi & 0\\
0 & \sin\varphi & \cos\varphi & 0\\
0 & 0 & 0 & 1
\end{pmatrix}
\qquad \text{(base } \{\ket{00},\ket{01},\ket{10},\ket{11}\}\text{)},
\label{eq:rbs}
\end{equation}
identità su $\ket{00}$ e $\ket{11}$, rotazione reale ordinaria sul piano
$\{\ket{01},\ket{10}\}$: per costruzione un operatore ortogonale che ruota
un solo piano coordinato, lasciando invariato il resto dello spazio.
Realizzato in porte elementari con due CZ attorno a una coppia di
rotazioni $R_y(\varphi),R_y(-\varphi)$ (Fig.~\ref{fig:rbs}): due porte a
due qubit.

\begin{figure}[htbp]
\centering
\includegraphics[width=0.55\textwidth]{figure/circ_rbs_espanso}
\caption{Blocco RBS$(\varphi)$ in porte elementari: due CZ attorno a due
rotazioni $R_y(\pm\varphi)$, tra Hadamard.}
\label{fig:rbs}
\end{figure}
\FloatBarrier

La scelta di una rotazione reale (RBS), invece dell'alternativa di
letteratura $W_{ij}(\theta)=e^{-i\theta\,\vs i\cdot\vs j}$ --- equivalente
in popolazioni ($\cos^2\varphi,\sin^2\varphi$ con $\theta=2\varphi$) ma con
una fase relativa immaginaria fra le due ampiezze, e una realizzazione a
tre porte a due qubit ($R_{xx}R_{yy}R_{zz}$) contro le due di RBS --- è
giustificata dal teorema spettrale: una matrice reale simmetrica come $H$
ammette sempre una base di autovettori reali, quindi un operatore che deve
portare uno stato reale in un altro stato reale, dentro un sottospazio a
due dimensioni, non ha bisogno di fasi complesse. La fase extra di $W$ è un
grado di libertà sempre sprecato per un'Hamiltoniana reale come questa, a
un costo maggiore in porte a due qubit --- costo che pesa ulteriormente
nel sistema aperto, dove ogni porta a due qubit è un canale di rumore in
più.

L'ansatz esteso aggiunge, dopo il blocco RBS, due rotazioni $R_y$
indipendenti sui due qubit (Fig.~\ref{fig:pma-esteso}): due parametri in
più (tre in totale), che rompono deliberatamente la conservazione di $M$
facendo uscire lo stato dal settore. A differenza dell'ansatz base, qui
non serve alcuna preparazione già entangled né una scelta di ramo in base
a $b/J$: la porta fissa iniziale è sempre una singola $X$ (stato di
prodotto $\ket{10}$), perché RBS esplora da sola l'intero sottospazio
$\{\ket{01},\ket{10}\}$ --- il singoletto stesso è uno dei punti che
raggiunge, a $\varphi=\pi/4$. L'entanglement, in questo ansatz, non è un
dato di partenza costruito da porte fisse: è il risultato della stessa
porta che l'ottimizzatore addestra.

\begin{figure}[htbp]
\centering
\includegraphics[width=0.42\textwidth]{figure/circ_pma_esteso}
\caption{Ansatz esteso: blocco RBS seguito da due rotazioni $R_y$
indipendenti sui due qubit.}
\label{fig:pma-esteso}
\end{figure}
\FloatBarrier

Con l'ansatz esteso la fidelity torna a $1$ su tutto il range di campo, con
e senza termine DM (Fig.~\ref{fig:riparazione}).

\begin{figure}[htbp]
\centering
\includegraphics[width=0.7\textwidth]{figure/fig04_vqe_riparazione}
\caption{Fidelity dell'ansatz esteso (RBS $+\,2$ rotazioni) rispetto allo
stato fondamentale esatto, con termine DM acceso: fidelity massima su
tutto il range.}
\label{fig:riparazione}
\end{figure}
\FloatBarrier

\subsection{Risultati numerici di riferimento}

Oltre alla scansione in campo, un punto di lavoro singolo --- indicato nel
seguito come punto \emph{P1}, $b/J=0.35$, $D/J=0.80$ --- è usato come
riferimento nelle sezioni successive, scelto perché lo stato fondamentale
vi è una sovrapposizione genuina fra settori e il termine DM è ben acceso:

\begin{table}[htbp]
\centering
\renewcommand{\arraystretch}{1.25}
\begin{tabular}{@{}lccc@{}}
\toprule
Ansatz & Parametri & $\mathcal F$ (punto P1) & CNOT (base nativa) \\
\midrule
HA (generico) & 6 & $1.000000$ & --- \\
PMA base & 1 & $0.940200$ & --- \\
PMA esteso (RBS $+\,2$) & 3 & $1.000000$ (12 cifre) & $2$ \\
\bottomrule
\end{tabular}
\caption{Fidelity al punto di lavoro P1 ($b/J=0.35$, $D/J=0.80$), usato
come riferimento nelle Sez.~\ref{sec:correlatori} e~\ref{sec:vqe-rumore}
in avanti. L'ansatz esteso, transpilato nella base nativa
$\{rz,sx,x,\mathrm{cx}\}$, costa $2$ CNOT.}
\label{tab:vqe-p1}
\end{table}

Un secondo punto di lavoro --- indicato come \emph{P2},
$b/J=-0.18$, $D/J=1$ (termine DM cinque volte più forte) --- è usato nelle
sezioni sulla dinamica temporale (Sez.~\ref{sec:trotter}), dove serve una
dinamica non monocromatica anziché uno stato fondamentale non banale. Il
confronto fra ansatz si conferma identico anche a questo punto: fidelity
HA $=1$, PMA base $=0.851664$, PMA esteso $=1$ (dodici cifre).

Il criterio di progetto che emerge da questa sezione è generale: un
ansatz variazionale va costruito sulle simmetrie che l'Hamiltoniana ha
davvero, non su quelle che sarebbe comodo avesse. Un ansatz che rispetta
una simmetria assente nell'Hamiltoniana è strutturalmente incapace di
raggiungere lo stato fondamentale, indipendentemente dal numero di
parametri; un blocco che rompe volutamente quella simmetria, anche con
pochi parametri aggiuntivi, la ripara.

\section{Evoluzione temporale: decomposizione di Suzuki--Trotter}
\label{sec:trotter}

\subsection{Il problema e la scomposizione dell'Hamiltoniana}

L'evoluzione temporale di uno stato è $\ket{\psi(t)}=e^{-iHt}\ket{\psi(0)}$
($\hbar=1$ in tutto il seguito). Non esiste una porta elementare che
realizzi direttamente $e^{-iHt}$ per un'Hamiltoniana arbitraria; la via
d'uscita è scomporre $H$ in pezzi che si sappiano esponenziare
singolarmente. Per il dimero,
\begin{equation}
H = \underbrace{J\,\vs1\!\cdot\!\vs2 + b\,(Z_1{+}Z_2)}_{H_1}
  \;+\; \underbrace{D\,(X_1Z_2-Z_1X_2)}_{H_2},
\label{eq:H1H2}
\end{equation}
dove $H_1$ raccoglie scambio isotropo e campo (i termini che, con $D=0,$
conservano la magnetizzazione totale) e $H_2$ è il solo termine DM. Con
$D=0$, $H_1$ e $H_2$ commutano banalmente ($H_2=0$) e non serve alcuna
approssimazione. Con $D\neq0$, $H_1$ e $H_2$ non commutano, quindi
$e^{-iHt}\neq e^{-iH_1t}e^{-iH_2t}$, e serve un metodo approssimato
controllato: la decomposizione di Suzuki--Trotter.

\subsection{Formula di Trotter ed errore}

Si divide l'intervallo $[0,t]$ in $N$ passi di durata $\tau=t/N$ e si
alterna
\begin{equation}
e^{-iHt} \;\simeq\; \Big(e^{-iH_2\tau}\,e^{-iH_1\tau}\Big)^{N}.
\label{eq:trotter}
\end{equation}
L'errore per singolo passo è $O(\tau^2)$ (dal commutatore
$[H_1,H_2]\neq0$); sommato su $N$ passi, l'errore sull'operatore di
evoluzione resta $O(t^2/N)$, cioè $O(1/N)$ a $t$ fissato, mentre la
perdita di fidelity sullo stato --- quadratica nell'errore
sull'operatore --- scala come $O(1/N^2)$. All'aumentare di $N$
l'approssimazione migliora, ma il circuito si allunga proporzionalmente.

Il circuito per un singolo passo (Fig.~\ref{fig:circ-trotter}) realizza
$e^{-iH_1\tau}$ con due rotazioni $R_z$ (per il campo) seguite dai tre
blocchi $R_{xx},R_{yy},R_{zz}$ (per lo scambio isotropo, i cui addendi
commutano a due a due e si fattorizzano quindi esattamente), e
$e^{-iH_2\tau}$ con due rotazioni $ZZ$ coniugate da porte di Hadamard.

\begin{figure}[htbp]
\centering
\includegraphics[width=\textwidth]{figure/circ_trotter}
\caption{Un singolo passo di Trotter: $e^{-iH_1\tau}$ (rotazioni $R_z$ del
campo, poi $R_{xx},R_{yy},R_{zz}$ dello scambio) seguito da
$e^{-iH_2\tau}$ (termine DM, due $R_{zz}$ coniugate da Hadamard). Il
circuito completo è questo blocco ripetuto $N$ volte.}
\label{fig:circ-trotter}
\end{figure}
\FloatBarrier

\subsection{Verifica dello scaling dell'errore}

Al punto di lavoro P2 ($b/J=-0.18$, $D/J=1$), con stato iniziale
$\ket{\uparrow\uparrow}$ (scelto perché non è autostato di $H_1$ né di
$H_2$: uno stato che lo fosse evolverebbe banalmente sotto ciascun blocco e
non mostrerebbe l'errore di Trotter anche quando presente), l'evoluzione
calcolata algebricamente e quella eseguita dal circuito transpilato
coincidono a precisione di macchina (scarto massimo $2.4\times10^{-13}$ su
una griglia da $N=1$ a $N=512$): l'implementazione in porte non introduce
errori propri oltre a quello di Trotter stesso.

\begin{figure}[htbp]
\centering
\includegraphics[width=0.72\textwidth]{figure/fig06_trotter_errore}
\caption{Errore a tempo fissato ($t=6$) al crescere di $N$, scala
doppio-logaritmica: distanza fra operatori di evoluzione (blu, scaling
$1/N$) e perdita di fidelity sullo stato (rosso, scaling $1/N^2$). Le
rette tratteggiate sono le pendenze teoriche attese.}
\label{fig:errore-trotter}
\end{figure}
\FloatBarrier

Le pendenze misurate con un fit sui punti $N\ge8$ sono $-1.008$ e
$-2.059$, contro i valori teorici $-1$ e $-2$ (Tab.~\ref{tab:errore-trotter}).

\begin{table}[htbp]
\centering
\renewcommand{\arraystretch}{1.2}
\begin{tabular}{@{}rccc@{}}
\toprule
$N$ & $\lVert U_{\rm circuito}-U_{\rm esatto}\rVert$ & $1-\mathcal F$ & $\lvert\Delta\media{S_z}\rvert$ \\
\midrule
2 & $1.26$ & $3.9\times10^{-1}$ & $1.8\times10^{-1}$ \\
8 & $2.7\times10^{-1}$ & $8.6\times10^{-3}$ & $3.8\times10^{-2}$ \\
32 & $6.5\times10^{-2}$ & $4.3\times10^{-4}$ & $2.3\times10^{-3}$ \\
128 & $1.6\times10^{-2}$ & $2.6\times10^{-5}$ & $1.5\times10^{-4}$ \\
512 & $4.0\times10^{-3}$ & $1.6\times10^{-6}$ & $9.1\times10^{-6}$ \\
\bottomrule
\end{tabular}
\caption{Errore a $t=6$, punto P2. Raddoppiando $N$, la prima colonna si
dimezza e la seconda si divide per quattro, come atteso dallo scaling
teorico.}
\label{tab:errore-trotter}
\end{table}

\subsection{Costo in porte}

Il passo di Trotter, scritto come in Fig.~\ref{fig:circ-trotter}, usa $5$
porte a due qubit. Ricompilando l'intero passo, invece di tradurre le $5$
porte una per una, nella base nativa (CNOT più rotazioni a un qubit), il
compilatore riconosce che insieme realizzano un unico operatore unitario a
due qubit e lo risintetizza con la decomposizione ottima nota per un gate
a due qubit arbitrario, al più $3$ CNOT:

\begin{table}[htbp]
\centering
\renewcommand{\arraystretch}{1.25}
\begin{tabular}{@{}lcc@{}}
\toprule
Compilazione & CNOT per passo & Profondità \\
\midrule
Traduzione diretta, non ottimizzata & $10$ & $37$ \\
Compilazione ottimizzata & $\mathbf{3}$ & $13$ \\
\bottomrule
\end{tabular}
\caption{Costo di un passo di Trotter dopo transpilazione nel set di porte
nativo. La compilazione ottimizzata riduce le porte a due qubit di un
fattore superiore a tre.}
\label{tab:costo-trotter}
\end{table}
\FloatBarrier

L'errore di Trotter scende aumentando $N$, ma il costo in porte a due
qubit --- il fattore dominante il tasso di errore su hardware reale ---
sale linearmente con $N$: esiste quindi un compromesso, e un $N$ ottimo
che non è né il più piccolo né il più grande. Determinarlo richiede un
modello di rumore esplicito (Sez.~\ref{sec:dinamica-rumore}).

\subsection{Combinazione con l'errore di preparazione}

Lo stato di partenza dell'evoluzione temporale, nella pipeline completa,
non è $\ket{\uparrow\uparrow}$ ma lo stato prodotto dal VQE
(Sez.~\ref{sec:vqe}). Ripetendo l'evoluzione di Trotter al punto P2 a
partire dai due ansatz --- la cui fidelity rispetto al fondamentale esatto
è, a questo punto, $0.8517$ per il PMA base e $1$ (dodici cifre) per il
PMA esteso --- i due errori del progetto (preparazione ed evoluzione) si
separano nettamente (Fig.~\ref{fig:trotter-due-ansatz}).

\begin{figure}[htbp]
\centering
\includegraphics[width=\textwidth]{figure/fig12_trotter_due_ansatz}
\caption{Evoluzione di $\media{S_z^{\rm tot}}(t)$ al punto P2, con $N=20$
passi di Trotter, a partire dai due ansatz del VQE. Nera: evoluzione
esatta dal fondamentale esatto. Tratteggiata: evoluzione esatta a partire
dallo stato del VQE (isola l'errore di preparazione). Punteggiata:
evoluzione con il circuito di Trotter a partire dallo stesso stato
(entrambi gli errori). \textbf{Sinistra}, PMA base: le curve del VQE si
staccano nettamente dall'esatta. \textbf{Destra}, PMA esteso: le tre
curve sono indistinguibili.}
\label{fig:trotter-due-ansatz}
\end{figure}
\FloatBarrier

\begin{table}[htbp]
\centering
\renewcommand{\arraystretch}{1.3}
\begin{tabular}{@{}lccc@{}}
\toprule
Ansatz & Solo VQE & VQE + circuito ($N=20$) & Solo circuito (da VQE) \\
\midrule
PMA base (1 par.) & $0.327$ & $0.322$ & $0.024$ \\
PMA esteso (3 par.) & $0.000$ & $0.033$ & $0.033$ \\
\bottomrule
\end{tabular}
\caption{Scarto massimo su $\media{S_z^{\rm tot}}(t)$ al punto P2,
scomposto per contributo. Con l'ansatz base, l'errore di preparazione
($0.327$) domina di oltre un ordine di grandezza quello di Trotter a
$N=20$ ($0.024$): aumentare $N$ non lo tocca. Con l'ansatz esteso l'errore
di preparazione si annulla per costruzione e resta solo Trotter
($0.033$).}
\label{tab:doppio-errore}
\end{table}

I due errori del progetto non sono intercambiabili: conviene sempre
attaccare per primo l'errore dominante, prima di investire in un $N$ più
grande.

\section{Correlatori dinamici a due tempi}
\label{sec:correlatori}

\subsection{Definizione e problema di misura}

La grandezza di interesse è il correlatore a due tempi
\begin{equation}
C_{ij}^{\alpha\beta}(t) \;=\;
\bra{\psi_0}\,\sigma_i^\alpha(t)\,\sigma_j^\beta(0)\,\ket{\psi_0},
\qquad \alpha,\beta\in\{x,y,z\},\quad i,j\in\{1,2\},
\label{eq:correlatore}
\end{equation}
la cui trasformata di Fourier è la quantità osservata in una misura di
scattering di neutroni su una molecola magnetica. Il prodotto
$\sigma_i^\alpha(t)\sigma_j^\beta(0)$ non è hermitiano: il suo valore
d'aspettazione è in generale complesso, e nessuna misura proiettiva lo
restituisce direttamente.

\subsection{Circuito: test di Hadamard con qubit ausiliario}

Il correlatore si misura con un qubit ausiliario (ancilla) in
sovrapposizione, che condiziona le operazioni sul registro fisico
(Fig.~\ref{fig:circ-correlatore}): si prepara $\ket{\psi_0}$ sul registro,
si apre l'ancilla con Hadamard, si applica $\sigma_j^\beta$ controllata
dall'ancilla, si evolve il registro con $U(t)$ --- \emph{non} controllata,
condizione essenziale per la realizzabilità del circuito --- si applica
$\sigma_i^\alpha$ anti-controllata, si richiude l'ancilla e la si misura.
Parte reale e parte immaginaria del correlatore si ottengono da due
esecuzioni che differiscono solo per l'ultima rotazione sull'ancilla
prima della misura ($H$ per la parte reale, $R_x(\pi/2)$ per
l'immaginaria).

\begin{figure}[htbp]
\centering
\includegraphics[width=0.78\textwidth]{figure/circ_correlatore}
\caption{Circuito ad ancilla per il correlatore dinamico (schema con soli
$2$ passi di Trotter per leggibilità). L'evoluzione $U(t)$ non è
controllata; solo gli operatori di Pauli lo sono.}
\label{fig:circ-correlatore}
\end{figure}
\FloatBarrier

Verificato contro un riferimento classico indipendente (esponenziale di
matrice diretto, che non condivide codice con il circuito): residuo medio
$4.5\times10^{-3}$, massimo $1.1\times10^{-2}$ con $N=200$ passi di
Trotter, crescente con $t$ come atteso da un errore di discretizzazione
(Fig.~\ref{fig:validazione-correlatore}).

\begin{figure}[htbp]
\centering
\includegraphics[width=0.8\textwidth]{figure/fig07_correlatore_validazione}
\caption{Correlatore rappresentativo su $t\in[0,6]$: curve continue
(riferimento classico) contro simboli (circuito quantistico), parte reale
e immaginaria.}
\label{fig:validazione-correlatore}
\end{figure}
\FloatBarrier

\subsection{Scelta del correlatore e contenuto spettrale}

Con due siti e tre componenti esistono $36$ combinazioni
$(i,\alpha,j,\beta)$, non equivalenti. Dodici sono esattamente nulle
all'istante iniziale --- le combinazioni a siti diversi con una sola
componente $y$, e quelle sullo stesso sito con componenti $x,z$ --- perché
l'Hamiltoniana è reale e simmetrica, il suo stato fondamentale si può
scegliere reale, e un prodotto di Pauli hermitiano ma puramente
immaginario ha media nulla su uno stato reale. Nessuna delle $36$ resta
nulla per ogni tempo, tuttavia: uno zero iniziale non implica un segnale
trascurabile a tempi successivi (Fig.~\ref{fig:scan36}).

\begin{figure}[htbp]
\centering
\includegraphics[width=\textwidth]{figure/fig08_scan36}
\caption{Tutte le $36$ combinazioni al punto P1. \textbf{Sinistra:} valore
all'istante iniziale, dodici caselle esattamente nulle. \textbf{Destra:}
ampiezza massima su $t\in[0,8]$, nessuna casella nulla.}
\label{fig:scan36}
\end{figure}
\FloatBarrier

Inserendo la risoluzione dell'identità sugli autostati $\ket n$ di $H$,
\begin{equation}
C_{ij}^{\alpha\beta}(t) = \sum_{n} c_n\, e^{-i(E_n-E_0)t}, \qquad
c_n = \bra0\sigma_i^\alpha\ket n\bra n\sigma_j^\beta\ket0,
\label{eq:decomposizione-spettrale}
\end{equation}
il segnale è esattamente una somma di quattro esponenziali complessi (uno
per autostato del dimero), verificata numericamente a precisione di
macchina contro il calcolo diretto. Un correlatore è "ricco" (più
transizioni di peso comparabile, rapporto $a_2/a_1$ fra le due righe più
pesanti $\approx1$) o "monocromatico" ($a_2/a_1\approx0$) indipendentemente
dalla sua ampiezza: le due proprietà non sono correlate. Il correlatore
$C_{11}^{yz}$, con $a_2/a_1=0.71$ (tre pesi comparabili), è scelto come
riferimento nel seguito, incluso nella Parte 2, perché genuinamente ricco
senza essere un caso estremo.

\subsection{Effetto di una preparazione imperfetta}

Sostituendo le ampiezze esatte con lo stato prodotto dal VQE al punto P1
($\mathcal F=0.9402$ per il PMA base, $\mathcal F\approx1$ per il PMA
esteso) e usando il circuito completo (ansatz $+$ ancilla $+$ $N=200$
passi di Trotter) per l'evoluzione, si osserva la stessa gerarchia di
errori già vista in Sez.~\ref{sec:trotter} (Fig.~\ref{fig:correlatore-vqe}).

\begin{figure}[htbp]
\centering
\includegraphics[width=\textwidth]{figure/fig14_correlatore_da_vqe}
\caption{$\lvert C_{11}^{yz}(t)\rvert$ su $t\in[0,8]$, preparato con i due
ansatz e fatto evolvere con il circuito completo. \textbf{Sinistra}, PMA
base: scarto visibile dal riferimento esatto. \textbf{Destra}, PMA
esteso: curve pressoché indistinguibili.}
\label{fig:correlatore-vqe}
\end{figure}
\FloatBarrier

Lo scarto massimo sul correlatore complesso, circuito vero incluso, è
$0.607$ con l'ansatz base --- dominato quasi per intero dall'errore di
preparazione, che da solo vale $0.602$ --- e $0.019$ con l'ansatz esteso,
dove l'errore di preparazione è nullo per costruzione e resta solo
l'errore di Trotter. Uno scan sistematico sulle $36$ combinazioni, con
entrambi gli ansatz e il circuito vero, conferma il quadro: per il PMA
base uno scarto medio $\sim0.35$ con un fattore $\sim11\times$ fra
combinazione più protetta e più fragile (le direzioni legate a una
rotazione $y$, la stessa usata dal blocco dell'ansatz base, restano
relativamente protette); per il PMA esteso uno scarto uniformemente
piccolo, $0.011$--$0.032$ --- il vero errore di Trotter residuo, non più
mascherato da un errore di preparazione dominante.

\section{Sistema aperto: il modello di rumore}
\label{sec:rumore}

\subsection{Dal vettore di stato all'operatore densità}

Tutti i risultati fin qui discussi assumono un'evoluzione unitaria su uno
stato puro $\ket\psi$: un'idealizzazione utile per mettere a punto un
metodo, ma non ciò che avviene su un dispositivo quantistico reale, dove
ogni qubit è accoppiato a un ambiente non osservato e ogni gate ha una
probabilità non nulla di eseguire un'operazione diversa da quella
prevista. L'Hamiltoniana \eqref{eq:H} non cambia; cambia il modo in cui il
circuito viene eseguito. Lo stato del sistema smette di essere un vettore
puro e diventa un operatore densità $\rho$; l'evoluzione, non più
un'unitaria pura, è un'operazione quantistica più generale (in Qiskit,
`AerSimulator(method="density\_matrix")`). Ogni quantità della pipeline si
riscrive di conseguenza: l'energia diventa $\Tr[\rho H]$ al posto di
$\bra\psi H\ket\psi$, e la fedeltà rispetto a un riferimento puro
$\ket{\psi_{\rm rif}}$ diventa $F=\bra{\psi_{\rm rif}}\rho\ket{\psi_{\rm rif}}$.

\subsection{Due canali di rumore, calibrati su hardware reale}

Il modello di rumore adottato è deliberatamente minimale: due soli
meccanismi, l'errore di gate (canale depolarizzante, separato a uno e due
qubit) e l'errore di lettura (readout). Rilassamento e defasamento
($T_1$, $T_2$) sono esclusi in questa modellazione.

Il canale depolarizzante a $n$ qubit ($d=2^n$) sostituisce lo stato con
una miscela fra lo stato originale e lo stato massimamente misto,
\begin{equation}
\mathcal{E}(\rho) = (1-\lambda)\,\rho + \lambda\,\frac{\mathbf 1}{d}.
\label{eq:depolarizzante}
\end{equation}
Le calibrazioni pubblicate riportano l'errore medio di gate $\varepsilon$
da randomized benchmarking, legato a $\lambda$ da
$\varepsilon=\lambda(d-1)/d$, da cui $\lambda_{1q}=2\varepsilon_{1q}$ e
$\lambda_{2q}=\tfrac43\varepsilon_{2q}$. Il readout simmetrico è descritto
da un solo parametro $p$ (probabilità di lettura errata, uguale in
entrambe le direzioni), che agisce come fattore puramente moltiplicativo
sull'osservabile $\media Z$,
\begin{equation}
\media{Z}_{\rm readout} = (1-2p)\,\media{Z}_{\rm ideale}.
\label{eq:readout-simmetrico}
\end{equation}
Questa proprietà --- apparentemente un dettaglio --- è decisiva più avanti
(Sez.~\ref{sec:scan-rumore}): un fattore moltiplicativo positivo non
sposta mai la posizione di un massimo.

\begin{figure}[htbp]
\centering
\includegraphics[width=0.55\textwidth]{figure/fig_depolarizzazione_bloch}
\caption{Contrazione del vettore di Bloch, $1-\lambda$, in funzione di
$\lambda$. Ai valori di calibrazione di riferimento la contrazione è
piccola (sotto l'1\%) ma non nulla.}
\end{figure}
\FloatBarrier

I quattro numeri del modello sono le mediane di calibrazione pubblicate
per \texttt{ibm\_torino} (IBM Heron r1, 133 qubit), da Stenger, Bazargan,
Bronn e Gunlycke~\cite{stenger2025}:

\begin{table}[htbp]
\centering
\renewcommand{\arraystretch}{1.3}
\begin{tabular}{@{}llc@{}}
\toprule
Canale & Gate di riferimento & Valore mediano\\
\midrule
Errore di gate, 1 qubit & \texttt{sx} & $\varepsilon_{1q}=2.9\times10^{-4}$\\
Errore di gate, 2 qubit & \texttt{CZ} & $\varepsilon_{2q}=3.8\times10^{-3}$\\
Errore di lettura & --- & $p_{\rm readout}=2.3\times10^{-2}$\\
\bottomrule
\end{tabular}
\caption{Parametri di calibrazione usati nel modello di rumore.}
\label{tab:calibrazione}
\end{table}

Il gate nativo a due qubit di \texttt{ibm\_torino} è il CZ, mentre il
circuito di questa tesi è scritto in una base che include il CNOT.
L'identità circuitale esatta
$\text{CNOT}=(\mathbf 1\otimes H)\cdot\text{CZ}\cdot(\mathbf 1\otimes H)$
(verificata anche numericamente a precisione di macchina) mostra che
l'interazione fisica a due qubit --- dove si concentra la maggior parte
dell'errore --- è la stessa in entrambi i casi; ciò che cambia è un cambio
di base realizzato da due gate a un qubit, il cui costo aggiuntivo
($\sim2\varepsilon_{1q}$, circa il 15\% di $\varepsilon_{2q}$) resta
dichiarato. È quindi corretto descrivere $\varepsilon_{2q}$(CZ) come valore
rappresentativo dell'ordine di grandezza dell'errore a due qubit, non come
una replica esatta del CNOT sintetizzato.

Con questi quattro numeri, il \texttt{NoiseModel} risultante coinvolge
esattamente le istruzioni $\{rz,sx,x,\mathrm{cx},\mathrm{measure}\}$: i
gate della base $\{rz,sx,x\}$ portano il canale a un qubit, il
\texttt{cx} quello a due qubit, la \texttt{measure} l'errore di lettura.
Ogni canale è applicato indipendentemente qubit per qubit (nessun
crosstalk); il readout è simmetrico in questo modello di base (l'ipotesi
è rilassata in Sez.~\ref{sec:readout-asimmetrico}); la calibrazione è un
solo punto nel tempo.

\section{Stato fondamentale sotto rumore}
\label{sec:vqe-rumore}

Con il modello di rumore fissato, la prima domanda è la più semplice
della catena: l'ansatz esteso del VQE (Sez.~\ref{sec:vqe}), con gli stessi
parametri $\theta^*$ ottimizzati nel caso ideale e riusati senza
rioptimizzare, trova ancora uno stato accettabile se il circuito gira su
un simulatore rumoroso? (La domanda complementare --- rioptimizzare
vedendo il rumore --- è trattata in Sez.~\ref{sec:noise-aware}.) Il
circuito, transpilato nella base nativa $\{rz,sx,x,\mathrm{cx}\}$, costa
$2$ CNOT, al punto di lavoro P1.

\begin{table}[htbp]
\centering
\renewcommand{\arraystretch}{1.3}
\begin{tabular}{@{}lccc@{}}
\toprule
Condizione & $E$ & $F$ & CNOT\\
\midrule
Rumore nullo (coerenza) & $-3.57321145$ & $1.0000000000$ & $2$\\
Rumore di riferimento (\texttt{ibm\_torino}) & $-3.50164605$ & $0.9850060074$ & $2$\\
\bottomrule
\end{tabular}
\caption{Energia e fedeltà del VQE al punto P1, senza e con rumore di
riferimento. Il caso a rumore nullo riproduce l'energia esatta a
precisione numerica.}
\label{tab:vqe-rumore}
\end{table}

Al rumore di riferimento la perdita di fedeltà è $1-F=0.014994$, poco
sotto l'1.5\%: la scala di degradazione più piccola di tutta la pipeline,
perché il circuito di preparazione, a soli $2$ CNOT, lascia poco spazio
al rumore di accumularsi. Al variare di $\varepsilon_{2q}$ la degradazione
è monotona e liscia, senza soglie (Fig.~\ref{fig:vqe-scan-eps2q}).

\begin{figure}[htbp]
\centering
\includegraphics[width=0.85\textwidth]{figure/fig_passo2_scan_eps2q}
\caption{Energia e fedeltà rumorose in funzione di $\varepsilon_{2q}$, a
$\varepsilon_{1q}$ e $p_{\rm readout}$ fissati al valore di riferimento
(linea tratteggiata). Degradazione monotona, nessuna soglia.}
\label{fig:vqe-scan-eps2q}
\end{figure}
\FloatBarrier

\section{Dinamica sotto rumore: il numero ottimale di passi}
\label{sec:dinamica-rumore}

\subsection{Il compromesso}

Nel sistema chiuso, aumentare $N$ migliora sempre l'approssimazione
dell'evoluzione. Sotto rumore questo non è più vero: ogni passo aggiuntivo
introduce gate rumorosi, e deve quindi esistere un numero ottimale di
passi $N^*$. La quantità di riferimento è la fedeltà dello stato rumoroso
rispetto allo stato bersaglio fisico,
\begin{equation}
\ket{\psi_{\rm target}(t)} = e^{-iHt}\ket{\psi_0^{\rm esatto}}, \qquad
F(N) = \bra{\psi_{\rm target}(t)}\rho_N\ket{\psi_{\rm target}(t)},
\label{eq:F-N}
\end{equation}
dove $\rho_N$ è ottenuto da: preparazione VQE rumorosa (Sez.~\ref{sec:vqe-rumore})
seguita da $N$ passi di Trotter rumorosi. Il bersaglio
$\ket{\psi_{\rm target}(t)}$ è calcolato dall'esponenziale diretto
dell'Hamiltoniana intera (DM incluso), non dalla scomposizione di Trotter:
un confronto altrimenti circolare. Una singola quantità impacchetta tre
contributi: la preparazione (costante in $N$), l'errore di Trotter
(decrescente in $N$) e il rumore di gate accumulato (crescente in $N$).

Il singolo passo di Trotter è transpilato \emph{una sola volta} al suo
costo minimo (matrice $4\times4$ risintetizzata nella base nativa, al più
$3$ CNOT, lo stesso limite già stabilito in Sez.~\ref{sec:trotter}), e il
blocco già compilato è composto $N$ volte; transpilare invece l'intero
circuito assemblato produrrebbe una fusione spuria degli $N$ passi
identici in un unico unitario a costo costante, cancellando la dipendenza
da $N$ che è l'oggetto dello studio.

\subsection{Risultato}

\begin{table}[htbp]
\centering
\renewcommand{\arraystretch}{1.15}
\small
\begin{tabular}{@{}cccc@{}}
\toprule
$N$ & $F$ (rumore nullo) & $F$ (\texttt{ibm\_torino}) & \\
\midrule
$1$ & $0.7799372$ & $0.7567291$ & \\
$2$ & $0.4776280$ & $0.4624084$ & \\
$3$ & $0.4018000$ & $0.3892000$ & minimo non perturbativo\\
$8$ & $0.9454850$ & $0.8177249$ & $\approx N^*$\\
$20$ & $0.9921835$ & $0.7097565$ & \\
$80$ & $0.9995212$ & $0.3567462$ & \\
$160$ & $0.9998804$ & $0.2686760$ & \\
\bottomrule
\end{tabular}
\caption{Fedeltà rispetto allo stato bersaglio fisico al punto P1, senza
e con rumore di riferimento. Senza rumore $F(N)\to1$ monotonamente per
$N\gtrsim5$; con rumore la curva cresce fino a $N=8$ e poi decresce.}
\label{tab:F-N}
\end{table}

\begin{figure}[htbp]
\centering
\includegraphics[width=0.68\textwidth]{figure/fig_passo3_F_vs_N}
\caption{Fedeltà $F(N)$ rispetto allo stato bersaglio fisico, a rumore di
riferimento. Crescita rapida per $N$ piccolo (l'errore di Trotter
domina), massimo a $N^*=8$, decrescita lenta per $N$ grande (il rumore di
gate accumulato domina).}
\label{fig:F-N}
\end{figure}
\FloatBarrier

\begin{center}
\renewcommand{\arraystretch}{1.2}
\begin{tabular}{@{}lc@{}}
\toprule
Quantità & Valore\\
\midrule
$N^*$ & $8$\\
$F(N^*)$ & $0.817725$\\
\bottomrule
\end{tabular}
\end{center}

Solo il rumore di gate accumulato crea il massimo: senza rumore l'ottimo
sarebbe $N\to\infty$. Un minimo non perturbativo a $N=3$ ($F\approx0.40$),
identico a rumore nullo e a rumore acceso, conferma che si tratta di un
effetto di Trotter puro nel regime non ancora perturbativo ($\tau=t/N$
non piccolo per $N=1,2,3$), non di un artefatto del rumore o della
simulazione. Il readout non compare in questa metrica --- è una fedeltà
fra stati, non una misura via ancilla --- ed entra per la prima volta nel
correlatore dinamico.

\section{Correlatori sotto rumore}
\label{sec:correlatori-rumore}

\subsection{Il circuito, con il rumore su ogni ingrediente}

Lo schema è lo stesso della Sez.~\ref{sec:correlatori}: preparazione di
$\ket{\psi_0}$, Hadamard sull'ancilla, controlled-$W$, evoluzione $U(t)$
via Trotter (non controllata), anti-controlled-$V$, rotazione di base e
misura sull'ancilla. Sotto rumore, ciascun ingrediente è imperfetto: la
preparazione (Sez.~\ref{sec:vqe-rumore}), il blocco $U(t)$ ripetuto
(Sez.~\ref{sec:dinamica-rumore}), e la misura finale, che include anche
l'errore di lettura --- qui il readout entra per la prima volta come una
vera istruzione di misura, con un \texttt{ReadoutError} genuino agganciato
al simulatore e $\media Z$ stimato dai conteggi su $2\times10^5$ shot
(incertezza statistica attesa $\sigma\approx0.0022$), mai da una formula
applicata a posteriori.

Il conteggio di gate cresce linearmente in $N$: $4+3N$ CNOT (preparazione
$2$, controlled-$W$ $1$, anti-controlled-$V$/rotazione $1$, più $3$ per
ogni passo di Trotter), verificato blocco per blocco.

Il correlatore di riferimento resta $C_{11}^{yz}$, scelto in
Sez.~\ref{sec:correlatori} per il suo contenuto spettrale ricco.

\subsection{Risultato: $N^*$ del correlatore}

\begin{table}[htbp]
\centering
\renewcommand{\arraystretch}{1.15}
\small
\begin{tabular}{@{}cccc@{}}
\toprule
$N$ & $|C|$ (rumore nullo) & $|C|$ (\texttt{ibm\_torino}, shot finiti) & \\
\midrule
$1$ & $0.4844$ & $0.4339$ & \\
$2$ & $0.6574$ & $0.5894$ & $N^*$\\
$3$ & $0.4998$ & $0.4298$ & \\
$6$ & $0.2181$ & $0.1733$ & minimo locale\\
$20$ & $0.2753$ & $0.1630$ & \\
\bottomrule
\end{tabular}
\caption{Modulo del correlatore $C_{11}^{yz}(t{=}2)$ in funzione di $N$,
punto P1.}
\label{tab:C-N}
\end{table}

\begin{figure}[htbp]
\centering
\includegraphics[width=0.68\textwidth]{figure/fig_passo4_correlatore_vs_N}
\caption{$\lvert C_{11}^{yz}(t{=}2,N)\rvert$ in funzione di $N$, readout
genuinamente campionato. Massimo a $N^*=2$, minimo locale attorno a
$N\approx6$, lenta risalita successiva.}
\label{fig:C-N}
\end{figure}
\FloatBarrier

A differenza della fedeltà di Trotter (Sez.~\ref{sec:dinamica-rumore}),
qui il massimo non nasce da un compromesso graduale Trotter/rumore: già a
rumore nullo $|C(N)|$ ha un massimo locale non perturbativo a $N=2$
($0.657$, quasi il doppio del valore fisico esatto
$|C_{\rm esatto}(t{=}2)|=0.318$, ottenuto dalla somma spettrale
\eqref{eq:decomposizione-spettrale}), seguito da un minimo attorno a
$N\approx6$ e da una lenta risalita verso il valore vero
(Fig.~\ref{fig:C-N-confronto}). Il rumore di gate accentua un massimo già
presente nella dinamica di Trotter pura; non lo crea.

\begin{figure}[htbp]
\centering
\includegraphics[width=0.68\textwidth]{figure/fig_passo4_confronto_rumore}
\caption{$|C_{11}^{yz}(t{=}2,N)|$ senza e con rumore, $N=1$--$40$: il
picco a $N=2$, il minimo a $N\approx6$ e la lenta risalita successiva sono
presenti in entrambe le curve.}
\label{fig:C-N-confronto}
\end{figure}
\FloatBarrier

Va sottolineato che $N^*=\arg\max_N|C(N)|$ individua il punto di segnale
più misurabile sotto rumore, non la stima più accurata del correlatore
fisico $C(t)$: qui i due criteri divergono di un fattore $\sim2$. Per gli
scopi di questa sezione e della successiva (fissare un compromesso
Trotter/rumore, verificarne la sensibilità ai parametri del canale) è il
criterio corretto; una stima fisicamente accurata di $C(t)$ andrebbe
corretta per questo effetto.

\subsection{Rappresentatività: scan sulle 36 combinazioni}

Ripetendo lo scan sistematico di Sez.~\ref{sec:correlatori} sotto rumore,
con il circuito vero e readout campionato ($5\times10^4$ shot per punto,
incertezza attesa $\approx0.0045$), la distribuzione di $N^*$ sulle $36$
combinazioni è concentrata su valori piccoli:

\begin{center}
\renewcommand{\arraystretch}{1.2}
\begin{tabular}{@{}cc@{}}
\toprule
$N^*$ & combinazioni (su 36)\\
\midrule
$1$ & $14$\\
$2$ & $12$\\
$3$ & $6$\\
$4$ & $2$\\
$5$ & $2$\\
\bottomrule
\end{tabular}
\end{center}

$C_{11}^{yz}$ cade nel secondo gruppo più popoloso ($N^*=2$): un caso
rappresentativo della distribuzione, non un estremo. Il fattore che
decide $N^*$ non è la ricchezza spettrale $a_2/a_1$ ma quanto guadagno non
perturbativo la dinamica di Trotter offre appena oltre $N=1$: le
combinazioni che partono già vicine al proprio valore asintotico
collassano a $N^*=1$ (nessun compromesso da negoziare); quelle con un
guadagno sostanziale fra $N=1$ e il primo picco locale (dall'$11\%$ al
$1229\%$, verificato su tutte le $36$ curve a rumore nullo) giustificano
$N^*$ più alti.

\section{Dipendenza di $N^*$ dai parametri di rumore}
\label{sec:scan-rumore}

Il compromesso Trotter/rumore misurato al punto di calibrazione di
riferimento non è universale: dipende dal livello di rumore. Esplorando
$\varepsilon_{1q}$ e $\varepsilon_{2q}$ un parametro alla volta (l'altro
fissato al valore di riferimento):

\begin{center}
\renewcommand{\arraystretch}{1.15}
\small
\begin{tabular}{@{}cc@{}}
\toprule
$\varepsilon_{2q}$ & $N^*$\\
\midrule
$1.0\times10^{-3}$ & $10$\\
$3.8\times10^{-3}$ (riferimento) & $8$\\
$6.0\times10^{-3}$ & $7$\\
$1.0\times10^{-2}$ & $1$\\
$2.5\times10^{-2}$ & $1$\\
\bottomrule
\end{tabular}
\qquad
\begin{tabular}{@{}cc@{}}
\toprule
$\varepsilon_{1q}$ & $N^*$\\
\midrule
$1.0\times10^{-4}$ & $9$\\
$2.9\times10^{-4}$ (riferimento) & $8$\\
$1.0\times10^{-3}$ & $7$\\
$3.0\times10^{-3}$ & $1$\\
\bottomrule
\end{tabular}
\end{center}

\begin{figure}[htbp]
\centering
\includegraphics[width=0.9\textwidth]{figure/fig_passo5_Nstar_vs_eps}
\caption{$N^*$ (fedeltà di Trotter) in funzione di $\varepsilon_{2q}$
(sinistra) e $\varepsilon_{1q}$ (destra), scala log. Monotonia
non-crescente, con transizione a $N^*=1$ per rumore sufficientemente
forte.}
\label{fig:Nstar-vs-eps}
\end{figure}
\FloatBarrier

$N^*$ è monotono non-crescente in entrambi i parametri di gate: oltre una
soglia di rumore, nessun numero di passi maggiore di uno vale il costo
aggiuntivo in gate rumorosi. Un modello ingenuo per l'errore totale,
$\varepsilon_{\rm tot}(N)\sim a/N+g\varepsilon_{2q}N$ (Trotter più rumore,
entrambi additivi), predice una pendenza log-log $-0.5$ per
$N^*(\varepsilon_{2q})$; quella osservata sui punti non saturati è
$-0.195$ --- compatibile con un decadimento moltiplicativo tipico di una
fedeltà accumulata su rumore depolarizzante, $F(N)\sim(1-\varepsilon_{2q})^{cN}$,
piuttosto che con un errore additivo su una singola quantità.

Il readout, per contro, non sposta la posizione di $N^*$: verificato su
misure genuinamente campionate, $N^*=2$ per il correlatore $C_{11}^{yz}$
per ogni $p_{\rm readout}$ testato dallo $0\%$ al $20\%$ (ben oltre il
valore di calibrazione, $2.3\%$), come atteso dal fatto che il readout
simmetrico agisce come fattore puramente moltiplicativo
\eqref{eq:readout-simmetrico}: un riscalamento per una costante positiva
non sposta mai la posizione di un argmax.

\section{Estensioni: ottimizzazione noise-aware e readout asimmetrico}
\label{sec:estensioni}

\subsection{VQE noise-aware}
\label{sec:noise-aware}

I risultati delle Sez.~\ref{sec:vqe-rumore}--\ref{sec:scan-rumore}
misurano la degradazione di parametri VQE ottimizzati \emph{senza} vedere
il rumore. Resta da chiedersi se rioptimizzare vedendo il rumore già
durante l'ottimizzazione porti a parametri sostanzialmente diversi, e se
questo cambi le conclusioni a valle.

Il canale depolarizzante è isotropo: per lo stato puro
$\ket{\psi(\theta)}$ preparato dall'ansatz,
\begin{equation}
\rho(\theta) = (1-\lambda)\,\ket{\psi(\theta)}\bra{\psi(\theta)} + \lambda\,\frac{\mathbf 1}{d},
\end{equation}
sicché l'energia rumorosa è una trasformazione \emph{affine}
dell'energia ideale,
\begin{equation}
E_{\rm rumoroso}(\theta) = (1-\lambda)\,E_{\rm ideale}(\theta) + \lambda\,\frac{\Tr[H]}{d}.
\label{eq:affine-energia}
\end{equation}
Una trasformazione affine a coefficiente positivo non sposta la posizione
del minimo: lo stesso $\theta^*$ che minimizza $E_{\rm ideale}$ minimizza
anche $E_{\rm rumoroso}$. A ciò si aggiunge che il conteggio CNOT del
circuito transpilato non dipende dai valori numerici di $\theta$
(verificato su assegnazioni casuali indipendenti), quindi non c'è nemmeno
una via indiretta, tramite il costo in gate, per cui convenga uno
spostamento.

La previsione --- $\theta^*_{\rm rumore}\approx\theta^*_{\rm ideale}$,
nessun vantaggio misurabile a rioptimizzare --- è confermata numericamente
su due metriche strutturalmente diverse (una fedeltà additiva, il modulo
di un correlatore), due punti di lavoro fisici (P1 e P2) e l'intera
griglia di parametri di rumore della Sez.~\ref{sec:scan-rumore}
($25/25$ punti):

\begin{table}[htbp]
\centering
\renewcommand{\arraystretch}{1.25}
\small
\begin{tabular}{@{}lcc@{}}
\toprule
Quantità & Riuso $\theta^*_{\rm ideale}$ & VQE noise-aware\\
\midrule
$E$ (rumoroso, P1) & $-3.50164605$ & $-3.50164616$\\
$F$ (rumorosa, P1) & $0.98500601$ & $0.98500602$\\
$N^*$ fedeltà di Trotter (P1) & $8$ & $8$\\
$N^*$ correlatore $C_{11}^{yz}$ (P1) & $2$ & $2$\\
\bottomrule
\end{tabular}
\caption{Confronto fra il riuso dei parametri ideali e la rioptimizzazione
noise-aware. Lo scostamento in energia è $\Delta E\sim10^{-7}$, sei ordini
di grandezza sotto la scala fisica del problema ($1-F\sim1.5\times10^{-2}$);
$N^*$ resta identico su entrambe le metriche.}
\label{tab:noise-aware}
\end{table}

\begin{figure}[htbp]
\centering
\includegraphics[width=0.85\textwidth]{figure/fig_F_vs_N_confronto}
\caption{Fedeltà di Trotter $F(N)$, riuso ideale contro noise-aware:
massimo a $N^*=8$ in entrambi i casi; lo scostamento (pannello destro,
scala ingrandita) resta $\sim10^{-4}$.}
\end{figure}

\begin{figure}[htbp]
\centering
\includegraphics[width=0.85\textwidth]{figure/fig_correlatore_noise_aware_confronto}
\caption{Modulo del correlatore rumoroso, riuso ideale contro
noise-aware: massimo a $N^*=2$ in entrambi i casi.}
\end{figure}
\FloatBarrier

Ripetuto a un secondo punto di lavoro (P2, termine DM cinque volte più
forte), $N^*$ resta invariato fra le due preparazioni anche lì
($8\to8$ per la fedeltà di Trotter, $2\to2$ per il correlatore). La
ragione generale non è specifica della struttura di $N^*$: per qualunque
osservabile ragionevolmente liscio dello stato preparato, uno scostamento
$\Delta\theta$ piccolo garantisce, per continuità, uno scostamento a
valle altrettanto piccolo, additivo o no.

\subsection{Readout asimmetrico}
\label{sec:readout-asimmetrico}

Il modello di readout usato fin qui è simmetrico, $p_{01}=p_{10}=p$. Un
readout asimmetrico ha due parametri distinti,
$p_{01}=P(\text{misuro }1\mid\text{stato }\ket0)$ e
$p_{10}=P(\text{misuro }0\mid\text{stato }\ket1)$; ripetendo la
derivazione della correzione su $\media Z$ con una matrice di confusione
generica,
\begin{equation}
\media{Z}_{\rm readout} = \alpha\,\media{Z}_{\rm ideale} + \beta,
\qquad \alpha \equiv 1-p_{01}-p_{10}, \qquad \beta \equiv p_{10}-p_{01}.
\label{eq:readout-asimmetrico}
\end{equation}
Il caso noto \eqref{eq:readout-simmetrico} è il sottocaso $\beta=0$ di
questa formula più generale. Fisicamente, il meccanismo dominante
dell'asimmetria è il rilassamento $T_1$ durante la finestra di misura: un
qubit in $\ket1$ può decadere a $\ket0$ prima che la misura sia completata,
mentre il processo inverso violerebbe la conservazione dell'energia in
assenza di eccitazione esterna --- da cui un'asimmetria strutturale
$p_{10}\gtrsim p_{01}$, osservata sistematicamente su hardware
superconduttore IBM, tipicamente con un fattore $2$--$4\times$.

Per il correlatore, la cui grandezza rilevante è un modulo nel piano
complesso, un offset costante $\beta\neq0$ trasla l'intera curva $C(N)$
nel piano complesso invece di riscalarla soltanto: la distanza
dall'origine di ciascun punto cambia in modo diverso a seconda della sua
posizione, e il massimo dopo la traslazione può non essere più quello di
prima. L'invarianza di $N^*$ dimostrata nel caso simmetrico non si estende
automaticamente: qui la risposta è empirica, non un teorema.

Nessuno split $(p_{01},p_{10})$ realmente registrato per
\texttt{ibm\_torino} è disponibile in forma citabile; si è quindi usato un
rapporto rappresentativo $p_{10}:p_{01}=3$, tratto da calibrazioni IBM
pubblicate su chip comparabili (range osservato $2$--$3.5\times$), a
parità di media con il valore di riferimento simmetrico
($p_{\rm readout}=2.3\times10^{-2}$): $p_{01}=0.0115$, $p_{10}=0.0345$.

\begin{table}[htbp]
\centering
\renewcommand{\arraystretch}{1.2}
\begin{tabular}{@{}lcc@{}}
\toprule
Configurazione di readout & $N^*$ & $|C(N^*)|$\\
\midrule
Simmetrico ($p_{01}=p_{10}=0.023$) & $2$ & $0.5869$\\
Asimmetrico illustrativo ($p_{01}=0.0115$, $p_{10}=0.0345$) & $2$ & $0.6195$\\
\bottomrule
\end{tabular}
\caption{Modulo del correlatore rumoroso al suo massimo, readout
simmetrico contro asimmetrico (stessa media), verificato su dati
genuinamente campionati.}
\label{tab:readout-asimmetrico}
\end{table}

\begin{figure}[htbp]
\centering
\includegraphics[width=0.68\textwidth]{figure/fig_correlatore_N_asimmetrico}
\caption{$\lvert C_{11}^{yz}(t{=}2,N)\rvert$, readout simmetrico contro
asimmetrico a parità di media: il massimo resta a $N^*=2$ in entrambi i
casi.}
\end{figure}
\FloatBarrier

Uno stress test a media fissa, con rapporti $p_{10}{:}p_{01}$ ben oltre il
range realistico, mostra $N^*=2$ invariato fino a un rapporto $50\times$:

\begin{center}
\renewcommand{\arraystretch}{1.15}
\small
\begin{tabular}{@{}cccc@{}}
\toprule
Rapporto $p_{10}{:}p_{01}$ & $p_{01}$ & $p_{10}$ & $N^*$\\
\midrule
$1$ (simmetrico) & $0.0230$ & $0.0230$ & $2$\\
$3$ (illustrativo) & $0.0115$ & $0.0345$ & $2$\\
$10$ & $0.0042$ & $0.0418$ & $2$\\
$50$ & $0.0009$ & $0.0451$ & $2$\\
\bottomrule
\end{tabular}
\end{center}

\begin{figure}[htbp]
\centering
\includegraphics[width=0.58\textwidth]{figure/fig_Nstar_vs_rapporto}
\caption{$N^*$ in funzione del rapporto $p_{10}{:}p_{01}$ (scala log), a
media fissa: $N^*=2$ per quattro ordini di grandezza di rapporto.}
\end{figure}
\FloatBarrier

Questa invarianza empirica non è garantita dalla teoria --- che ammette in
linea di principio anche il caso opposto, un offset che appiattisce o
sposta il massimo --- ma è coerente con essa: per questo correlatore, a
questo punto di lavoro, l'offset $\beta(1+i)$ introdotto dall'asimmetria
si allontana dall'origine più agli $N$ già vicini al picco che agli $N$
lontani da esso, accentuando il massimo esistente invece di spostarlo. È
un fatto numerico specifico di questo correlatore e di questo punto di
lavoro, non una conseguenza necessaria della struttura del problema.

\section{Sintesi dei risultati}

Il dimero di Heisenberg isotropo con termine di Dzyaloshinskii--Moriya si
è dimostrato un banco di prova completo per la simulazione quantistica di
sistemi di spin molecolari, sufficientemente semplice da ammettere un
riferimento esatto a ogni stadio e sufficientemente ricco da contenere
tutti gli ingredienti fisici rilevanti: un incrocio di livelli protetto da
simmetria e la sua rottura esplicita, un ansatz variazionale che fallisce
per ragioni strutturali di simmetria e la sua riparazione minimale, la
necessità di un'approssimazione controllata (Trotter) per l'evoluzione
temporale non banale, un osservabile a due tempi che richiede un circuito
ausiliario, e infine il comportamento non monotono di un algoritmo di
simulazione temporale in presenza di rumore realistico.

I risultati quantitativi principali sono riassunti in Tab.~\ref{tab:sintesi}.

\begin{table}[htbp]
\centering
\renewcommand{\arraystretch}{1.35}
\small
\begin{tabular}{@{}p{4.6cm}p{10.5cm}@{}}
\toprule
Stadio & Risultato principale \\
\midrule
Sistema chiuso, spettro & Incrocio esatto a $B/J=2$ ($D=0$); gap minimo $0.565$ con $D/J=0.2$ \\
VQE, ansatz $M$-conservante & Fallisce vicino all'incrocio con DM ($\mathcal F\to0.70$); riparato da RBS $+\,2$ rotazioni ($\mathcal F=1$ sempre) \\
Dinamica (Trotter), sistema chiuso & Errore $\sim1/N$ (operatore), $\sim1/N^2$ (fidelity); costo minimo $3$ CNOT/passo \\
Correlatori, sistema chiuso & Errore di preparazione domina l'errore di Trotter quando l'ansatz non è esteso \\
VQE sotto rumore & $1-F\approx1.5\%$ al punto di calibrazione di riferimento \\
Dinamica sotto rumore & $N^*=8$ (fedeltà); solo il rumore di gate crea il massimo \\
Correlatori sotto rumore & $N^*=2$ (correlatore $C_{11}^{yz}$); rappresentativo su 36 combinazioni \\
Dipendenza dai parametri di rumore & $N^*$ monotono non-crescente in $\varepsilon_{1q},\varepsilon_{2q}$; invariante nel readout simmetrico \\
VQE noise-aware & Nessun vantaggio misurabile a rioptimizzare ($\Delta E\sim10^{-7}$) \\
Readout asimmetrico & $N^*=2$ invariato fino a rapporti $50\times$ oltre il range realistico \\
\bottomrule
\end{tabular}
\caption{Sintesi dei risultati principali, sistema chiuso e sistema
aperto, dimero di Heisenberg con termine DM.}
\label{tab:sintesi}
\end{table}

\begin{thebibliography}{9}

\bibitem{crippa2021}
D.~Crippa, R.~Pisoni, F.~Tacchino, A.~Chiesa, S.~Carretta, \emph{Simulation
of quantum many-body dynamics for a molecular spin system on a
superconducting quantum processor}, Magnetochemistry \textbf{7}, 117 (2021).

\bibitem{stenger2025}
J.~D.~Stenger, M.~Bazargan, B.~S.~Bronn, D.~Gunlycke, \emph{Method for
simulating open-system dynamics using mid-circuit measurements on a
quantum computer}, arXiv:2504.15187 (2025).

\end{thebibliography}

\end{document}

